%! Equivalent Representations of Trigonometric Functions — Unit 3, Lesson 3.12 — Exit Ticket (Answer Key)
\documentclass[10pt]{article}
\usepackage{precalculus-article}
\usepackage{precalculus-key}   % pulls in -boxes; do NOT also load -boxes
\newcommand{\TallMath}[1]{$\displaystyle #1\rule[-1.4em]{0pt}{3.2em}$}

\begin{document}

\pageheader{Unit 3, Lesson 3.12}{Exit Ticket --- Answer Key}
\namedateperiod

\vspace{0.12in}

\begin{notesbox}{Exit Ticket --- Answer independently}

\textbf{1.}\quad Use the double-angle identity to rewrite $\sin^2\theta$
in terms of $\cos 2\theta$.

\vspace{4pt}
We know $\cos 2\theta = 1 - 2\sin^2\theta$. Solve for $\sin^2\theta$:

\vspace{6pt}
$\sin^2\theta = $ \ans{$\dfrac{1 - \cos 2\theta}{2}$}

\vspace{0.16in}
\textbf{2.}\quad Use the difference identity to find the exact value of
$\cos\!\left(\dfrac{\pi}{12}\right)$.

\vspace{4pt}
\textit{(Hint:} $\dfrac{\pi}{12} = \dfrac{\pi}{3} - \dfrac{\pi}{4}$\textit{)}

\vspace{4pt}
$\cos\!\left(\dfrac{\pi}{3} - \dfrac{\pi}{4}\right)
= \cos\dfrac{\pi}{3}\cos\dfrac{\pi}{4} + \sin\dfrac{\pi}{3}\sin\dfrac{\pi}{4}$

\vspace{6pt}
$= $ \ans{$\dfrac{1}{2}$} $\cdot$ \ans{$\dfrac{\sqrt{2}}{2}$}
$+$ \ans{$\dfrac{\sqrt{3}}{2}$} $\cdot$ \ans{$\dfrac{\sqrt{2}}{2}$}

\vspace{6pt}
$= $ \ans{$\dfrac{\sqrt{2}}{4} + \dfrac{\sqrt{6}}{4} = \dfrac{\sqrt{6}+\sqrt{2}}{4}$}

\vspace{0.16in}
\textbf{3.}\quad Solve $\cos^2\theta - \sin^2\theta = 0$ for $\theta \in [0, 2\pi)$.

\vspace{4pt}
$\cos^2\theta - \sin^2\theta = \cos$ \ans{$2\theta$} $= 0$

\vspace{6pt}
$\cos 2\theta = 0 \;\Rightarrow\; 2\theta = $ \ans{$\dfrac{\pi}{2},\;\dfrac{3\pi}{2},\;\dfrac{5\pi}{2},\;\dfrac{7\pi}{2}$}

\vspace{6pt}
$\theta = $ \ans{$\dfrac{\pi}{4},\;\dfrac{3\pi}{4},\;\dfrac{5\pi}{4},\;\dfrac{7\pi}{4}$}

\begin{teachernote}
Q1 reinforces the half-angle rearrangement. Q2 gives exact-value practice with the difference identity.
Q3 is efficient because $\cos^2\theta - \sin^2\theta = \cos 2\theta$ is a double-angle form; students who
try Pythagorean substitution first should still arrive at the same answer but with more work.
\end{teachernote}

\end{notesbox}

\end{document}
